\section{问题三：给定预算的条件资源配置}
\subsection{决策变量、支持域与成本}
将第一问的 17 域实测优良配方质心固定为外生候选 $p^{\rm ref}$。根据题目允许以前问配比作固定输入的要求，优化变量取 $(N_B,D_B,Q_B)$；因 A4 与 B6 没有联合实验，不能把此计算称为 $N,D,Q,p$ 四变量联合最优。模型响应采用式(\ref{eq:m1})，限定在 $N_B\in[0.070542,11.965825]$、$D_B\in[10,299.893]$、$Q_B\in[0.4,1]$。连续 $Q_B$ 的结果是 B6 离散水平间插值。

总成本按照题给附录 B 写为
\begin{equation}
C_{\rm tot}=6ND+\eta ND L_{\rm ctx}
+D\,[g(Q_B)-g(Q_{B,0})]_+,\qquad
g(Q_B)=10^7e^{6Q_B},
\label{eq:cost}
\end{equation}
主情景取 $\eta=2\times10^{-4}$、$L_{\rm ctx}=2048$、$Q_{B,0}=0.4$。式中 $N,D$ 是原始参数与 token 数；三个成本项均为 FLOPs。$g$ 的指数、幂及对数备选形式均是题设情景而非市场报价。求解 $\min L_1$，约束 $C_{\rm tot}\leq C$ 及前述支持域；逐候选回代成本、边界和目标值，并检查角点的一阶必要条件。

\subsection{外生上下文的临界值}
由式(\ref{eq:cost})的基础训练与注意力两项相等，得到
\begin{equation}
6ND=\eta NDL_{\rm ctx}\quad\Longrightarrow\quad
L_{\rm ctx}^{\rm crit}=\frac{6}{\eta}=30000.
\end{equation}
C7 的最大位置容量记录覆盖 2048、4096、8192、32768 和 131072；这些是模型可支持窗口，并非训练时实际采用的序列长度。32768 情景略高于临界值，此时注意力代理与基础训练之比为 $\eta L_{\rm ctx}/6\approx1.092$。因此本文在 C7 可行容量上做情景敏感性，不把 2048 视为所有模型的实测上下文。
\subsection{主情景结果与边界解释}
\begin{table}[htbp]\centering\caption{指数质量成本及 2048 上下文下的条件配置}
\small\begin{tabular}{rrrrrl}\toprule
预算 FLOPs & $N_B^*$ & $D_B^*$ & $Q_B^*$ & 预测 Loss & 活跃边界\\\midrule
$10^{19}$ & 0.12898 & 10.000 & 0.55744 & 3.3064 & $D$ 下界\\
$10^{21}$ & 2.3231 & 53.152 & 1.0000 & 2.3802 & $Q$ 上界\\
$10^{22}$ & 5.6433 & 249.410 & 1.0000 & 2.1676 & $Q$ 上界\\
$10^{24}$ & 11.966 & 299.893 & 1.0000 & 2.1100 & $N,D,Q$ 上界\\\bottomrule
\end{tabular}\end{table}
\fig{problem3/figures/budget_path.pdf}{0.87}{预算路径与活跃约束；高预算平台由观测支持域截断。}

将结构性转移预先定义为最优解中 $Q_B$ 的下界、内点、上界状态或 $N_B,D_B$ 支持域边界的活跃集合改变。沿预算网格扫描并对活跃集变更作二分定位：指数成本主情景约在 $8.94\times10^{19}$ FLOPs 脱离 $D$ 下界，在 $1.54\times10^{20}$ FLOPs 使质量达到上界；约在 $1.34\times10^{22}$ 和 $2.42\times10^{22}$ FLOPs 先后触及 $D,N$ 上界。这些阈值是当前模型和支持域的数值转折，不是现实训练系统的普适算力门槛；连续成本份额的平滑变化也不称为结构转移。
$10^{19}$ FLOPs 的连续最优质量 0.55744 不属于 B6 的原生离散水平；在其实际水平上重新优化 $N,D$ 得 $Q_B=0.6$，模型预测 Loss 增加 0.00524。此预算的 $D_B=10$ 恰好触及下界，将数据下界向内提高到 20B 时预测 Loss 从 3.3064 升到 3.4199，表明低预算结果受边界设定控制。$10^{21}$ FLOPs 下，用题设幂/对数质量成本替换指数成本，预测最优数据量分别为 48.923B/68.765B；不能把其中一式单独当作真实报价。

在 $10^{24}$ FLOPs 情景，最优点触及三项支持上界，仍有 97.582\% 预算未用。该“最优”只在现有训练数据范围内成立；不能据此认为现实高预算不值得继续扩大模型或数据。成本参数及 B6 质量机制的不确定性没有被窄 Bootstrap 区间完整覆盖。第一问配比质心相对训练均值的 1M 响应面预测差为 $-0.00623$，配方行 Bootstrap 的 95\% 区间为 $[-0.01715,0.00316]$，并跨过零；该差不能加入表中的 B6 绝对 Loss。
